Si assuma l'operazionale ideale con VOUT in regime lineare. Si considerino i valori delle
resistenze pari a quelli nominali (in particolare R1=R2). Si assuma inoltre per semplicit\'a
Vg=0.

\paragraph{Richiesta (a)}
Si discuta lo stato di polarizzazione del diodo in corrispondenza di ciascuna semi-onda di VS
e si commenti l'andamento osservato in V+.

\paragraph{Richiesta (b)}
Per ciascuno dei due stati si determini la relazione tra VOUT e VS e si giustifichi l'andamento
osservato per VOUT al punto 1a. Di che tipo di circuito si tratta?

\paragraph{Richiesta (c)}
Si discuta il tipo di filtro operato da R4-C1 e le caratteristiche del segnale in VM.

\paragraph{Richiesta (d)}
Si calcoli il valore atteso per il rapporto tra il livello di continua di VM e l'ampiezza di VS e
lo si confronti con la pendenza della retta di best-fit ottenuta al punto 1d.

\paragraph{Richiesta (e)}
Si dia una stima del valore atteso percentuale del ripple di VM, ovvero del
rapporto tra la sua escursione ed il suo valor medio.
